% !TEX root = chapter-standalone.tex

\section{Actions}
\label{sec:actions}

\todotext{Expand this intro a bit to explain the conceptual idea of actions better, and mention the key term ``state space''. Maybe move the robot on a grid example up to this intro as a concrete motivating example. Also, maybe we can motivate/explain why it makes sense to abstractly separate the set of ``instructions'' from the functions that model the dynamics... why is this a good idea? what are the benefits? maybe a second example can also underscore this, maybe one from linear algebra?}

\todotext{Because of our semicolon infix notation for function composition, the `more natural' version of actions to prefer is right actions rather than left actions (normally it is the other way around). So maybe we should start with right actions and prefer/emphasize them in our exposition (while of course still mentioning left actions and explaining them...)}

\todotext{This section is to be rewritten, with the 'covariant' (could we also say coordinate free?) version of the definition of action as the main one. Then what is currently called a right action should be 'braided' and then called a left action (but as a remark, no longer as a definition).}

\todotext{Add a remark that explains the term/concept of 'contravariant' action (there is material on this in devel at the moment. If we use the term 'contravarian morphism' somewhere, be sure we define it.)}

\todotext{Throughout this chapter, maybe we should use S as a default letter for the 'state set' rather A?}



The idea of ``actions'' appears across mathematics and in many areas of application -- for example in the topic of dynamical systems in engineering, or in the topic of automata in computer science. 

Here is the central pattern captured by the concept of an action: on the one hand, a monoid (or some other algebraic gadget) is given which models time or models a sequence of events, and on the other hand that monoid ``acts'' on a set that models the possible states of a system. This action encodes how the system changes over time or changes in reaction to a sequence of events or inputs. Before fitting this idea into mathematical definitions, let us begin with two motivating examples. 

   
\subsection{Motivating example: a robot on a grid}
   
    Consider a robot moving on a two-dimensional grid made up of 16 squares, as illustrated in \cref{fig:robot_example}. To label each of the 16 squares we use simple coordinates: labels that consist of an ``$x_1$-direction'' (horizontal) and an ``$x_2$-direction" (vertical), with each dimension having possible entries $1$, $2$, $3$ or $4$. In other words, the set of 16 squares is labeled by the set
\begin{equation}
    \setX = \setX_1 \cartprod \setX_2
\end{equation}
with $\setX_1 = \makeset{1, 2, 3, 4}$ and $\setX_2 = \makeset{1, 2, 3, 4}$. 

    \begin{marginfigure}
        \centering
        \includesag{120_robotcategory}
        \caption{Robot on a grid example. }
        \label{fig:robot_example}
    \end{marginfigure}

    We consider the position coordinates of the robot to be its ``state''.  For example, in the figure, the state of the robot is that its position is $\tup{1,1} \setin \setX$. 
    
    
    Next, consider the following set of instructions that we may give the robot,
\begin{equation}
\prin=\{\uparrow, \downarrow,\rightarrow,\leftarrow\},
\end{equation}
where each symbol is a command to ``move one square over'' in the corresponding direction  -- unless the movement would lead the robot to leave the grid, in which case no movement is undertaken. 


We can model this situation using the following  ``dynamics function'': 
    \begin{equation}
        \begin{aligned}
            \prdyn\colon \prin \cartprod \prst                & \sto \prst \\
            \tupp{\uparrow, \tup{\prsteln{1},\prsteln{2}}}    & \mapsto \tupp{\prsteln{1}, \min(\prsteln{2}+1,4)}, \\
            \tupp{\downarrow, \tup{\prsteln{1},\prsteln{2}}}  & \mapsto \tupp{\prsteln{1}, \max(\prsteln{2}-1,1)}, \\
            \tupp{\rightarrow, \tup{\prsteln{1},\prsteln{2}}} & \mapsto \tupp{\min(\prsteln{1}+1,4), \prsteln{2}}, \\
            \tupp{\leftarrow, \tup{\prsteln{1},\prsteln{2}}}  & \mapsto \tupp{\max(\prsteln{1}-1,1), \prsteln{2}}, \\
        \end{aligned}
    \end{equation}
    where the ``$\min$'' and ``$\max$'' ensure that the robot stays in within the boundaries of the grid.
    
    \todotext{Rework the above example so that it fits better with the action definition as we now have it...}
    
\devel{

\subsection{Motivating example: a discrete-time dynamical system}

\todotext{What is a nice motivating example here??}

}

\subsection{Monoid actions}

\begin{ctdefinition}[Monoid action]
    \label{def:monoid-cov-action}
    An \maindef{action} of a \SY{monoid} $\monoidA = \tup{\monoidAset, \mtimesof\monoidA, \idmon_\monoidA}$ on a set $\stateset$ is:

    \constit
    \begin{enumerate}
        \item A \SY{monoid homomorphism}
              \begin{equation}
                  \label{eq:act-monoid}
                  \act \colon \monoidA \fto \Endof\stateset.
              \end{equation}
    \end{enumerate}
\end{ctdefinition}

Unpacked, the above definition means that we have a function
\begin{equation}
     \act \colon  \monoidAset \to (\stateset \to \stateset)
\end{equation}
that satisfies the following two conditions.
 
First, for all $\monelA, \monelB \setin \monoidAset$,
\begin{equation}
    \label{eq:act-to-endo-comp-compatibility}
    \act(\monelA \mtimesof\monoidA \monelB) = \act(\monelA) \mtimesof{\Endof\stateset} \act(\monelB), 
\end{equation}
which encodes ``compatibility of $\act$ with the composition operations in the source and the target'';

Second, 
\begin{equation}\label{eq:act-neutral-element-condit}
    \act(\idmon_\monoidA) = \catidat\stateset,
\end{equation}
which encodes ``compatibility of $\act$ with the neutral elements in the source and the target''. 

Both conditions may be unpacked even further, in terms of individual elements of $\stateset$. 
Recall that the composition operation of $\Endof\stateset$ is the composition of functions, $\mapa \mtimesof{\Endof\stateset} \mapb = \mapa \mthen \mapb$ for $\mapa, \mapb \setin \Endof\stateset$, in which the function written first is applied first (\cref{exa:endofunction-monoid}). 
Hence the first condition says
\begin{equation}
    \label{eq:act-to-endo-comp-compatibility-elementwise}
    \forall \monelA, \monelB \setin \monoidAset \colon \forall \statesetel \setin \stateset \colon \quad \act(\monelA \mtimesof\monoidA \monelB)(\statesetel) = \act(\monelB)\bigl(\act(\monelA)(\statesetel)\bigr).
\end{equation}
In words: acting by $\monelA \mtimesof\monoidA \monelB$ is the same as first acting by $\monelA$ and then acting by $\monelB$.
The second condition says
\begin{equation}
    \forall \statesetel \setin \stateset \colon \quad \act(\idmon_\monoidA)(\statesetel) = \statesetel.
\end{equation}

\todotext{Connect to examples now here...}

\begin{example}[Matrices acting on vectors]
    \label{exa:matrix-action}
    Fix a natural number $n$, and let $\reals^{\ntimesn}$ denote the set of $\ntimesn$ real matrices.
    Consider the \SY{monoid} $\monoidA = \tup{\reals^{\ntimesn}, \mtimesof\monoidA, \idmat}$ whose composition operation is defined, for all $\matmor{A}, \matmor{B} \setin \reals^{\ntimesn}$, by
    \begin{equation}
        \label{eq:matrix-action-composition}
        \matmor{A} \mtimesof\monoidA \matmor{B} \definedas \matmor{B} \matmor{A},
    \end{equation}
    where the right-hand side denotes the usual product of matrices, and whose \SY{neutral element} is the identity matrix $\idmat$.
    This is indeed a \SY{monoid}: associativity follows from the associativity of matrix multiplication, since
    \begin{equation}
        (\matmor{A} \mtimesof\monoidA \matmor{B}) \mtimesof\monoidA \matmor{C} = \matmor{C} (\matmor{B} \matmor{A}) = (\matmor{C} \matmor{B}) \matmor{A} = \matmor{A} \mtimesof\monoidA (\matmor{B} \mtimesof\monoidA \matmor{C}),
    \end{equation}
    and the neutrality laws hold because $\matmor{A} \idmat = \matmor{A} = \idmat \matmor{A}$.
    (In other words, $\monoidA$ is the opposite of the monoid of $\ntimesn$ matrices under the usual matrix multiplication; see \cref{exa:opposite-matrices}.)

    We read $\matmor{A} \mtimesof\monoidA \matmor{B}$ as ``first $\matmor{A}$, then $\matmor{B}$''.
    This matches what matrices do to vectors: to transform a vector $\vectob{v} \setin \reals^n$ first by $\matmor{A}$ and then by $\matmor{B}$, we compute $\matmor{B}(\matmor{A}\vectob{v}) = (\matmor{B}\matmor{A})\vectob{v}$.

    We define an action of $\monoidA$ on the set $\reals^n$ of real $n$-dimensional (column) vectors by
    \begin{equation}
        \label{eq:matrix-action}
        \act \colon \monoidA \fto \Endof{\reals^n},
    \end{equation}
    where, for each $\matmor{A} \setin \reals^{\ntimesn}$,
    \begin{equation}
        \defmapperiod{\act(\matmor{A})}{\reals^n}{\to}{\reals^n}{\vectob{v}}{\matmor{A}\vectob{v}}
    \end{equation}
    Let us check that $\act$ is a \SY{monoid morphism}.
    For all $\matmor{A}, \matmor{B} \setin \reals^{\ntimesn}$ and all $\vectob{v} \setin \reals^n$ we have
    \begin{equation}
        \act(\matmor{A} \mtimesof\monoidA \matmor{B})(\vectob{v})
        = (\matmor{B}\matmor{A})\vectob{v}
        = \matmor{B}(\matmor{A}\vectob{v})
        = \act(\matmor{B})\bigl(\act(\matmor{A})(\vectob{v})\bigr)
        = \bigl(\act(\matmor{A}) \mtimesof{\Endof{\reals^n}} \act(\matmor{B})\bigr)(\vectob{v}),
    \end{equation}
    which is condition \cref{eq:act-to-endo-comp-compatibility-elementwise}, and
    \begin{equation}
        \act(\idmat)(\vectob{v}) = \idmat \vectob{v} = \vectob{v} = \catidat{\reals^n}(\vectob{v}),
    \end{equation}
    which is condition \cref{eq:act-neutral-element-condit}.

    Note that the reversed order in \cref{eq:matrix-action-composition} is essential.
    If we had instead equipped $\reals^{\ntimesn}$ with the usual matrix multiplication, $\matmor{A} \cdot \matmor{B} = \matmor{A}\matmor{B}$, then the same assignment $\matmor{A} \mapsto (\vectob{v} \mapsto \matmor{A}\vectob{v})$ would satisfy
    \begin{equation}
        \act(\matmor{A}\matmor{B}) = \act(\matmor{B}) \mtimesof{\Endof{\reals^n}} \act(\matmor{A}),
    \end{equation}
    with the order reversed, and so, for $n \geq 2$, it would not be a \SY{monoid morphism} into $\Endof{\reals^n}$.
\end{example}

\begin{example}[Scalar multiplication]
Let $\monoidA = \tup{\reals, \cdot, 1}$ be the monoid of real numbers, with multiplication as composition operation and the number $1$ as neutral element. Consider the set $\stateset = \reals^3$ of real 3-dimensional vectors. An action 
\begin{equation}
\act \colon \reals \fto \Endof\reals^3
\end{equation}
is given by scalar multiplication: for each $\lambda \in \reals$, 
\begin{equation}
\defmapperiod{\act(\lambda)}{\reals^3}{\to}{\reals^3}{\begin{pmatrix} x_1 \\ x_2 \\ x_3 \end{pmatrix}}{\begin{pmatrix} x_1 \cdot \lambda \\ x_2 \cdot \lambda \\ x_3 \cdot \lambda \end{pmatrix}}
\end{equation}
\end{example}

\devel{
\begin{example}[robot on a grid]
\todotext{connect the intro to the definition}
\end{example}

\begin{example}[Plant growing]
\todotext{connect the example in the semigroups section that describes states of a growing plant}
\end{example}


\todotext{Make a remark/example somewhere that any monoid acts on itself...}

}





\subsection{Uncurrying actions}

Recall that (right) currying is the operation of applying the canonical isomorphism
\begin{equation}
    \curry \colon \setC^{\setA \cartprod \setB} \to \pars{\setC^\setB}^\setA
\end{equation}
which transforms any function of the type
\begin{equation}
    \setA \cartprod \setB \to \setC
\end{equation}
into one of the type
\begin{equation}
   \setA \to (\setB \to \setC).
\end{equation}

A monoid action involves a function of the type
\begin{equation}
     \act \colon  \monoidAset \to (\stateset \to \stateset).
\end{equation}
If we apply the inverse $\curry^{-1}$ (``uncurrying'') to this function, 
we obtain a function of the type
\begin{equation}\label{eq:uncurried-action}
 \makecartprod{\monoidAset, \stateset} \sto \stateset.
 \end{equation}
Throughout mathematics and applications, actions are often presented in this uncurried form. 
Since we can tell the difference between $\act$ and its uncurried version based on their respective domains and codomains, we will use the same name, ``$\act$'', for both versions.


Now, so far, we have just used $\curry^{-1}$ to produce the function \cref{eq:uncurried-action} from a monoid action. However, $\act \colon  \monoidAset \to (\stateset \to \stateset)$ was not just any function, it was a monoid homomorphism $\act \colon \monoidA \fto \Endof\stateset$. What does that condition translate to in the uncurried version $\act \colon  \makecartprod{\monoidAset, \stateset} \sto \stateset$? The answer is summarized in the following definition.

\begin{ctdefinition}[Monoid action, uncurried]
\label{def:uncurried-action}
   
    An \emph{uncurried action} of a \SY{monoid}~$\monoidA = \tup{\monoidAset, \mtimesof\monoidA, \idmon_\monoidA}$ on a set~\stateset is:
    
    \constit
    \begin{enumerate}
        \item a function
              \begin{equation}
                  \label{eq:ract-monoid}
                  \act \colon \makecartprod{\monoidAset, \stateset} \sto \stateset.
              \end{equation}
    \end{enumerate}
    \condit
    \begin{enumerate}
        \item for all~$\statesetel \setin \stateset$ and $\monelA,\monelB\setin\monoidAset$,
              \begin{equation}
                  \label{eq:monoiad-ract-1cond}
                  \act(\monelA \mtimesof\monoidA \monelB, \statesetel) = \act(\monelB, \act(\monelA, \statesetel));
              \end{equation}
        \item for all~$\statesetel \setin \stateset$,
              \begin{equation}
                  \label{eq:monoiad-ract-2cond}
                  \act(\idmon_\monoidA, \statesetel) = \statesetel.
              \end{equation}
    \end{enumerate}
\end{ctdefinition}

\begin{gradedexercise}[\exname{UncurriedAction}]
    \label{ex:UncurriedAction}
Throughout this exercise, fix a monoid $\monoidA = \tup{\monoidAset, \mtimesof\monoidA, \idmon_\monoidA}$ and~a set $\stateset$. 
\begin{enumerate}
\item Given a function
    \begin{equation}
    \act \colon \makecartprod{\monoidAset, \stateset} \sto \stateset
    \end{equation}
    that satisfies \cref{def:uncurried-action}, prove that the curried function
    \begin{equation}
    \curry(\act) \colon  \monoidAset \to (\stateset \to \stateset)
    \end{equation}
    satisfies conditions \cref{eq:act-to-endo-comp-compatibility} and \cref{eq:act-neutral-element-condit} and is hence a monoid action. 
    \item Given a monoid action 
    \begin{equation}
    \act \colon  \monoidAset \to (\stateset \to \stateset)
    \end{equation}
    prove that its uncurried version 
    \begin{equation}
    \curry^{-1}(\act) \colon  \makecartprod{\monoidAset, \stateset} \sto \stateset
    \end{equation}
    satisfies \cref{def:uncurried-action}. 
\end{enumerate}
\end{gradedexercise}
\solutionof{UncurriedAction}

\devel{

\todotext{Rewrite this part to discuss the possibility of writing uncurried action with $\monoidAset$ on the right...}

\begin{remark}
    Left actions are often denoted using infix notation of the kind
    \begin{equation}
        \makecartprod{\monoidAset, \stateset} \sto \stateset, \ \tup{\monela, \statesetel} \mapsto \monela \cdot \statesetel
    \end{equation}
    or even simply
    \begin{equation}
        \makecartprod{\monoidAset, \stateset} \sto \stateset, \ \tup{\monela, \statesetel} \mapsto \monela \statesetel
    \end{equation}
    (in regard to the latter, think of the notation often used for multiplication of matrices with vectors).
    If we write our notation ``$\monela \mtimesof\monoidA \monelb$'' as ``$\monela \star \monelb$'' (or even simply ``$\monela \monelb$''), then REF reads as
    \begin{equation}
        \monela \cdot (\monelb \cdot \statesetel) = (\monela \star \monelb) \cdot \statesetel,
    \end{equation}
    or simply
    \begin{equation}
        \monela (\monelb \statesetel) = (\monela  \monelb)  \statesetel.
    \end{equation}
    In other words, we can view REF as a kind of ``associative law'', but where two of the three elements come from $\monoidAset$, while one of the elements comes from $\stateset$.
\end{remark}
}



\subsection{Semigroup actions}



% There are two versions for actions: a ``covariant'' and a ``contravariant'' version, depending on the order of the transformations.
% \footnote{
%     Traditionally, people have used the names ``left'' and ``right'' to describe the two types.
%     However, those names are connected to a particular choice of notation.
%     We name the two versions by highlighting the properties of covariance and contravariance.
% }

%\linkvideo{spring2021-actions:semi-actions} % Semigroup actions


\begin{ctdefinition}[Semigroup action]
    \label{def:semigroup-cov-action}
    An \maindef{action} of a \SY{semigroup} $\sgrpA = \tup{\sgrpAset, \mtimesof\sgrpA}$ on a set $\stateset$ is:

    \constit
    \begin{enumerate}
        \item A \SY{semigroup} morphism
              \begin{equation}
                  \label{eq:act-semigroup}
                  \act \colon \sgrpA \fto \Endof\stateset.
              \end{equation}
    \end{enumerate}
\end{ctdefinition}

\begin{ctdefinition}[Semigroup action, uncurried]
    \label{def:semigroup-uncurried-action}
    An \emph{uncurried action} of a \SY{semigroup} $\sgrpA = \tup{\sgrpAset, \mtimesof\sgrpA}$ on a set $\stateset$ is:

    \constit
    \begin{enumerate}
        \item a function
              \begin{equation}
                  \label{eq:semigroup-uncurried-action}
                  \act \colon \makecartprod{\sgrpAset, \stateset} \sto \stateset.
              \end{equation}
    \end{enumerate}
    \condit
    \begin{enumerate}
        \item for all $\statesetel \setin \stateset$ and $\sgrpelA, \sgrpelB \setin \sgrpAset$,
              \begin{equation}
                  \label{eq:semigroup-uncurried-action-1cond}
                  \act(\sgrpelA \mtimesof\sgrpA \sgrpelB, \statesetel) = \act(\sgrpelB, \act(\sgrpelA, \statesetel)).
              \end{equation}
    \end{enumerate}
\end{ctdefinition}

% \begin{ctdefinition}
%     [Semigroup contravariant action (preliminary version)]\label{def:semigroup-contra-action-prelim}
%     A \emph{semigroup \textbf{contravariant} action} of a \SY{semigroup}~$\sgrpA$ on a set~$\stateset$ is a map
%     \begin{equation}
%         \label{eq:ract}
%         \ract \colon \makecartprod{\sgrpAset, \stateset} \to \stateset
%     \end{equation}
%     such that, for all~$\statesetel \setin \stateset$,
%     \begin{equation}
%         \label{eq:ract1cond}
%         \ract(\monelb, \ract(\monela, \statesetel)) = \ract(\monelb \mthenof\sgrpA  \monela, \statesetel).
%     \end{equation}
% \end{ctdefinition}


\devel{
\subsection{Actions of sets}

\begin{ctdefinition}
    [Set left action]\label{def:set-left-action-operation}
    A \emph{left action} of a set $\sgrpAset$ on a set $\stateset$ is:

    \constit
    \begin{enumerate}
        \item A function
              \begin{equation}
                  \label{eq:lact-set}
                  \act \colon \makecartprod{\sgrpAset, \stateset} \sto \stateset.
              \end{equation}
    \end{enumerate}
\end{ctdefinition}

\begin{ctdefinition}%
    [Set right action]\label{def:set-right-action-operation}
    A \emph{right action} of a set $\sgrpAset$ on a set $\stateset$ is:

    \constit
    \begin{enumerate}
        \item A function
              \begin{equation}
                  \label{eq:ract-set}
                  \act \colon \makecartprod{\stateset, \sgrpAset} \sto \stateset.
              \end{equation}
    \end{enumerate}
\end{ctdefinition}

\todotext{say how we can turn a set action into a semigroup action or a monoid action}

}



    \subsection{Contravariant actions}
    
    
    \begin{ctdefinition}[Contravariant monoid action]
        \label{def:monoid-cont-action}
        A \maindef{contravariant action} of a \SY{monoid} $\monoidA = \tup{\monoidAset, \mtimesof\monoidA, \idmon_\monoidA}$ on a set $\stateset$ is:

        \constit
        \begin{enumerate}
            \item A \SY{monoid morphism}
                  \begin{equation}
                      \label{eq:cont-act-monoid}
                      \act \colon \monoidA\op \fto \Endof\stateset.
                  \end{equation}
        \end{enumerate}
    \end{ctdefinition}

  

    % Diagram for the two conditions of def:monoid-uncurried-cont-action (corrected 2026-10-14), commented out for now:
    % \begin{marginfigure}
    %     \centering
    %     \includesag{semigroup-contrav-action}
    %     \caption{The conditions for an uncurried contravariant monoid action.}
    %     \label{fig:semigroup-contrav-action}
    % \end{marginfigure}
    
    
    \begin{ctdefinition}[Contravariant monoid action, uncurried]
        \label{def:monoid-uncurried-cont-action}
        An \emph{uncurried contravariant action} of a \SY{monoid} $\monoidA = \tup{\monoidAset, \mtimesof\monoidA, \idmon_\monoidA}$ on a set $\stateset$ is:

        \constit
        \begin{enumerate}
            \item a function
                  \begin{equation}
                      \label{eq:uncurried-cont-action}
                      \act \colon \makecartprod{\monoidAset, \stateset} \sto \stateset.
                  \end{equation}
        \end{enumerate}
        \condit
        \begin{enumerate}
            \item for all $\statesetel \setin \stateset$ and $\monelA, \monelB \setin \monoidAset$,
                  \begin{equation}
                      \label{eq:uncurried-cont-action-1cond}
                      \act(\monelA \mtimesof\monoidA \monelB, \statesetel) = \act(\monelA, \act(\monelB, \statesetel));
                  \end{equation}
            \item for all $\statesetel \setin \stateset$,
                  \begin{equation}
                      \label{eq:uncurried-cont-action-2cond}
                      \act(\idmon_\monoidA, \statesetel) = \statesetel.
                  \end{equation}
        \end{enumerate}
    \end{ctdefinition}

    \begin{remark}
        Just as in the covariant case (\cref{def:monoid-cov-action}, \cref{def:uncurried-action}, and \cref{ex:UncurriedAction}), one can prove that \cref{def:monoid-cont-action} and \cref{def:monoid-uncurried-cont-action} are equivalent: currying and uncurrying give a one-to-one correspondence between contravariant actions $\act \colon \monoidA\op \fto \Endof\stateset$ and uncurried contravariant actions $\act \colon \makecartprod{\monoidAset, \stateset} \sto \stateset$.
        The proof is the same as for \cref{ex:UncurriedAction}, with $\monoidA$ replaced by $\monoidA\op$: since $\monelA \mtimesof{\monoidA\op} \monelB = \monelB \mtimesof\monoidA \monelA$, the order of the two elements on the left-hand side of \cref{eq:monoiad-ract-1cond} is reversed, which gives \cref{eq:uncurried-cont-action-1cond}.
    \end{remark}



    In contrast to \cref{eq:monoiad-ract-1cond}, \cref{eq:uncurried-cont-action-1cond} says that, if we first act by $\monelB$ to obtain $\act(\monelB, \statesetel)$, and then act by $\monelA$, the result is the same as acting by $\monelA \mtimesof\monoidA \monelB$ (note the inverted order).

    % And what about contravariant actions?
    A \SY{contravariant semigroup action} can be defined as a covariant action of \emph{the opposite} \SY{semigroup}.

   




